% GENERATED FILE. Edit canonical lessons, then run npm run generate:books.
% canonical-source-hash: fnv1a64:b855115ebf117f1b
% canonical-lessons: UR-C345-qadr, UR-C345-iman, UR-C345-izzat, UR-C345-usul, UR-C345-aqida

\chapter{What We Value}
\label{ch:ur-what-we-value}
\hlchaptermodality{345}

\begin{chapteropening}
\textbf{By the end of this chapter:} \emph{I can say what a family or a neighbourhood values.}

Complete the last lesson of chapter 345: I can say what a family or a neighbourhood values.
\end{chapteropening}

\section[\texorpdfstring{\ur{قدر} (qadr)}{قدر (qadr)}]{\ur{قدر} (qadr) — worth; esteem}

\label{lesson:UR-C345-qadr}



\begin{quote}
Before the new one: say the Urdu for oppression, then the Urdu for to provoke.
\end{quote}

\subsection*{You'll want to know: \ur{قدر}}
\textbf{\ur{قدر}} — \emph{qadr} — \textquotedblleft{}worth; esteem\textquotedblright{}.

\begin{quote}
\ur{وہ} \ur{پیسے} \ur{کی} \ur{قدر} \ur{جانتا} \ur{ہے۔}
\end{quote}

\begin{quote}
\emph{voh paise kī qadr jāntā hai.}
\end{quote}

\begin{quote}
He knows the value of money.
\end{quote}

Feminine.

\begin{cousinweb}
\textbf{\ur{قدر}} (\emph{qadr}) is usually traced, through Persian, to Arabic \emph{qadr}, \textquotedblleft{}measure, worth\textquotedblright{}.
\end{cousinweb}

\begin{grammarlens}[title={What you've built}]
One more word for giving reasons and opinions.
\end{grammarlens}

\subsection*{Your turn}
\begin{itemize}
\raggedright
  \item \emph{Say it:} \emph{qadr}
  \item \emph{Say it:} \emph{qadr}, once more
  \item \emph{Say it:} \emph{voh paise kī qadr jāntā hai.} — He knows the value of money.
  \item \emph{Recall:} say the Urdu for oppression, then the Urdu for to provoke, then say \emph{qadr} again
\end{itemize}

\begin{tcolorbox}[breakable,colback=teal!4,colframe=teal!35!black,title={Before you move on}]
What does \ur{قدر} mean? (\textquotedblleft{}Worth; esteem\textquotedblright{}.) Say it once more.
\end{tcolorbox}
\section[\texorpdfstring{\ur{ایمان} (īmān)}{ایمان (īmān)}]{\ur{ایمان} (īmān) — faith; good faith}

\label{lesson:UR-C345-iman}



\begin{quote}
Before the new one: say the Urdu for to provoke, then the Urdu for worth.
\end{quote}

\subsection*{You'll want to know: \ur{ایمان}}
\textbf{\ur{ایمان}} — \emph{īmān} — \textquotedblleft{}faith; good faith\textquotedblright{}.

\begin{quote}
\ur{اس} \ur{کا} \ur{ایمان} \ur{مضبوط} \ur{ہے۔}
\end{quote}

\begin{quote}
\emph{us kā īmān mazbūt hai.}
\end{quote}

\begin{quote}
Her faith is strong.
\end{quote}

Masculine.

\begin{cousinweb}
\textbf{\ur{ایمان}} (\emph{īmān}) is usually traced, through Persian, to Arabic \emph{īmān}, \textquotedblleft{}faith\textquotedblright{}; \textbf{\ur{ایماندار}}, \textquotedblleft{}honest\textquotedblright{}, is built on it.
\end{cousinweb}

\begin{grammarlens}[title={What you've built}]
One more word for giving reasons and opinions.
\end{grammarlens}

\subsection*{Your turn}
\begin{itemize}
\raggedright
  \item \emph{Say it:} \emph{īmān}
  \item \emph{Say it:} \emph{īmān}, once more
  \item \emph{Say it:} \emph{us kā īmān mazbūt hai.} — Her faith is strong.
  \item \emph{Recall:} say the Urdu for to provoke, then the Urdu for worth, then say \emph{īmān} again
\end{itemize}

\begin{tcolorbox}[breakable,colback=teal!4,colframe=teal!35!black,title={Before you move on}]
What does \ur{ایمان} mean? (\textquotedblleft{}Faith; good faith\textquotedblright{}.) Say it once more.
\end{tcolorbox}
\section[\texorpdfstring{\ur{عزت} (izzat)}{عزت (izzat)}]{\ur{عزت} (izzat) — honour; respect}

\label{lesson:UR-C345-izzat}



\begin{quote}
Before the new one: say the Urdu for worth, then the Urdu for faith.
\end{quote}

\subsection*{You'll want to know: \ur{عزت}}
\textbf{\ur{عزت}} — \emph{izzat} — \textquotedblleft{}honour; respect\textquotedblright{}.

\begin{quote}
\ur{محلے} \ur{میں} \ur{اس} \ur{کی} \ur{بڑی} \ur{عزت} \ur{ہے۔}
\end{quote}

\begin{quote}
\emph{mahalle meṅ us kī baṛī izzat hai.}
\end{quote}

\begin{quote}
He is greatly respected in the neighbourhood.
\end{quote}

Feminine.

\begin{cousinweb}
\textbf{\ur{عزت}} (\emph{izzat}) is usually traced, through Persian, to Arabic \emph{ʿizza}, \textquotedblleft{}might, honour\textquotedblright{}.
\end{cousinweb}

\begin{grammarlens}[title={What you've built}]
One more word for giving reasons and opinions.
\end{grammarlens}

\subsection*{Your turn}
\begin{itemize}
\raggedright
  \item \emph{Say it:} \emph{izzat}
  \item \emph{Say it:} \emph{izzat}, once more
  \item \emph{Say it:} \emph{mahalle meṅ us kī baṛī izzat hai.} — He is greatly respected in the neighbourhood.
  \item \emph{Recall:} say the Urdu for worth, then the Urdu for faith, then say \emph{izzat} again
\end{itemize}

\begin{tcolorbox}[breakable,colback=teal!4,colframe=teal!35!black,title={Before you move on}]
What does \ur{عزت} mean? (\textquotedblleft{}Honour; respect\textquotedblright{}.) Say it once more.
\end{tcolorbox}
\section[\texorpdfstring{\ur{اصول} (usūl)}{اصول (usūl)}]{\ur{اصول} (usūl) — a principle}

\label{lesson:UR-C345-usul}



\begin{quote}
Before the new one: say the Urdu for faith, then the Urdu for honour.
\end{quote}

\subsection*{You'll want to know: \ur{اصول}}
\textbf{\ur{اصول}} — \emph{usūl} — \textquotedblleft{}a principle\textquotedblright{}.

\begin{quote}
\ur{میرے} \ur{کچھ} \ur{اصول} \ur{ہیں۔}
\end{quote}

\begin{quote}
\emph{mere kuchh usūl haiṅ.}
\end{quote}

\begin{quote}
I have some principles.
\end{quote}

Masculine.

\begin{cousinweb}
\textbf{\ur{اصول}} (\emph{usūl}) is usually traced, through Persian, to Arabic \emph{uṣūl}, \textquotedblleft{}roots, foundations\textquotedblright{}; Urdu uses the word as a singular, \textquotedblleft{}a principle\textquotedblright{}.
\end{cousinweb}

\begin{grammarlens}[title={What you've built}]
One more word for giving reasons and opinions.
\end{grammarlens}

\subsection*{Your turn}
\begin{itemize}
\raggedright
  \item \emph{Say it:} \emph{usūl}
  \item \emph{Say it:} \emph{usūl}, once more
  \item \emph{Say it:} \emph{mere kuchh usūl haiṅ.} — I have some principles.
  \item \emph{Recall:} say the Urdu for faith, then the Urdu for honour, then say \emph{usūl} again
\end{itemize}

\begin{tcolorbox}[breakable,colback=teal!4,colframe=teal!35!black,title={Before you move on}]
What does \ur{اصول} mean? (\textquotedblleft{}A principle\textquotedblright{}.) Say it once more.
\end{tcolorbox}
\section[\texorpdfstring{\ur{عقیدہ} (aqīdā)}{عقیدہ (aqīdā)}]{\ur{عقیدہ} (aqīdā) — a belief; a creed}

\label{lesson:UR-C345-aqida}



\begin{quote}
Before the new one: say the Urdu for honour, then the Urdu for a principle.
\end{quote}

\subsection*{You'll want to know: \ur{عقیدہ}}
\textbf{\ur{عقیدہ}} — \emph{aqīdā} — \textquotedblleft{}a belief; a creed\textquotedblright{}.

\begin{quote}
\ur{یہ} \ur{اس} \ur{کا} \ur{پرانا} \ur{عقیدہ} \ur{ہے۔}
\end{quote}

\begin{quote}
\emph{yih us kā purānā aqīdā hai.}
\end{quote}

\begin{quote}
This is an old belief of hers.
\end{quote}

Masculine.

\begin{cousinweb}
\textbf{\ur{عقیدہ}} (\emph{aqīdā}) is usually traced, through Persian, to Arabic \emph{ʿaqīda}, \textquotedblleft{}a creed, a firm belief\textquotedblright{}.
\end{cousinweb}

\begin{grammarlens}[title={What you've built}]
One more word for giving reasons and opinions.
\end{grammarlens}

\subsection*{Your turn}
\begin{itemize}
\raggedright
  \item \emph{Say it:} \emph{aqīdā}
  \item \emph{Say it:} \emph{aqīdā}, once more
  \item \emph{Say it:} \emph{yih us kā purānā aqīdā hai.} — This is an old belief of hers.
  \item \emph{Recall:} say the Urdu for honour, then the Urdu for a principle, then say \emph{aqīdā} again
\end{itemize}

\begin{tcolorbox}[breakable,colback=teal!4,colframe=teal!35!black,title={Before you move on}]
What does \ur{عقیدہ} mean? (\textquotedblleft{}A belief; a creed\textquotedblright{}.) Say it once more.
\end{tcolorbox}
